% !TEX program = xelatex
% ============================================================================
%  POSN Computer Qualification — Algorithms (Pseudo Code)
%  Topic: ขั้นตอนวิธี (Algorithms)
% ============================================================================

\documentclass[11pt, aspectratio=169]{beamer}
% ============================================================================
%  POSN Computer Qualification — Shared Beamer Preamble
%  Used by all topic slide decks. Compile with XeLaTeX.
% ============================================================================

% ---------- Thai language & font ----------
\usepackage{iftex}
\usepackage{fontspec}

% XeTeX-specific: enable Thai line breaking
\ifXeTeX
  \XeTeXlinebreaklocale{th}
\fi
% Noto Sans Thai — body text (clean modern sans-serif, handles all Thai)
\setmainfont{Noto Sans Thai}[
  UprightFont = Noto Sans Thai Light,
  BoldFont    = Noto Sans Thai Medium,
  Scale       = 0.9
]
\setsansfont{Noto Sans Thai}[
  UprightFont = Noto Sans Thai Light,
  BoldFont    = Noto Sans Thai Medium,
  Scale       = 0.9
]

% Noto Sans Thai — clean modern sans-serif for structural elements
\newfontfamily\headingfont{Noto Sans Thai}[
  UprightFont = Noto Sans Thai Medium,
  BoldFont    = Noto Sans Thai Bold,
  Scale       = 0.9
]

% --- Heading font for structural elements ---
\setbeamerfont{title}{family=\headingfont, series=\bfseries, size=\huge}
\setbeamerfont{subtitle}{family=\headingfont, series=\mdseries}
\setbeamerfont{author}{family=\headingfont, series=\mdseries}

% --- Custom Title Page ---
\defbeamertemplate{title page}{custom}[1][]{%
  \centering
  \vspace*{2cm}
  % Title — in white
  {\usebeamerfont{title}\color{white}\inserttitle\par}
  \vspace{0.6cm}
  % Subtitle — in white
  {\usebeamerfont{subtitle}\color{white}\insertsubtitle\par}
  \vspace{2.5cm}
  % Author — blue filled rounded box, white text
  \begin{center}
    \begin{tcolorbox}[arc=3mm, colback=boxBlue, colframe=boxBlue!80,
                      width=0.42\textwidth, halign=center,
                      left=1.5em, right=1.5em, top=0.8em, bottom=0.8em,
                      boxrule=1.5pt, coltext=white]
      \usebeamerfont{author}\insertauthor
    \end{tcolorbox}
  \end{center}
}
\setbeamertemplate{title page}[custom]
\setbeamerfont{frametitle}{family=\headingfont, series=\bfseries}
\setbeamerfont{section in toc}{family=\headingfont}
\setbeamerfont{page number in head/foot}{family=\headingfont}
\setbeamerfont{section title}{family=\headingfont, series=\bfseries}
\setbeamerfont{subsection title}{family=\headingfont}

% Monospace font (for code/verbatim)
\setmonofont{Latin Modern Mono}[Scale=0.9]

% ---------- Math ----------
\usepackage{amsmath, amssymb}

% ---------- Graphics ----------
\usepackage{graphicx}
\usepackage{tikz}

% ---------- String parsing (for section headers) ----------
\usepackage{xstring}

% ---------- Colored boxes ----------
\usepackage[skins, breakable, xparse]{tcolorbox}
% Note: we do NOT load the 'most' option — it predefines example/solution environments
% that would conflict with our custom boxes.

% --- Color definitions ---
\definecolor{boxBlue}{RGB}{41, 128, 185}
\definecolor{boxGreen}{RGB}{39, 174, 96}
\definecolor{boxOrange}{RGB}{230, 126, 34}
\definecolor{boxPurple}{RGB}{142, 68, 173}
\definecolor{boxBlueBg}{RGB}{235, 245, 255}
\definecolor{boxGreenBg}{RGB}{235, 255, 240}
\definecolor{boxOrangeBg}{RGB}{255, 245, 235}
\definecolor{boxPurpleBg}{RGB}{248, 240, 255}

% --- Explanation box (blue) — for definitions, concepts, notes ---
\newtcolorbox{explanation}[1][]{
  enhanced,
  breakable,
  arc         = 2mm,
  colback     = boxBlueBg,
  colframe    = boxBlue!50,
  title       = #1,
  fonttitle   = \bfseries\color{white},
  coltitle    = white,
  colbacktitle= boxBlue,
  attach boxed title to top left = {yshift=-2mm, xshift=2mm},
  boxed title style = {arc=2mm, size=small},
  before skip = 8pt,
  after skip  = 8pt
}

% --- Example box (green) — for worked examples ---
\newtcolorbox{exambox}[1][]{
  enhanced,
  breakable,
  arc         = 2mm,
  colback     = boxGreenBg,
  colframe    = boxGreen!50,
  title       = #1,
  fonttitle   = \bfseries\color{white},
  coltitle    = white,
  colbacktitle= boxGreen,
  attach boxed title to top left = {yshift=-2mm, xshift=2mm},
  boxed title style = {arc=2mm, size=small},
  before skip = 8pt,
  after skip  = 8pt
}

% --- Solution box (orange) — for solutions and answers ---
\newtcolorbox{solbox}[1][]{
  enhanced,
  breakable,
  arc         = 2mm,
  colback     = boxOrangeBg,
  colframe    = boxOrange!50,
  title       = #1,
  fonttitle   = \bfseries\color{white},
  coltitle    = white,
  colbacktitle= boxOrange,
  attach boxed title to top left = {yshift=-2mm, xshift=2mm},
  boxed title style = {arc=2mm, size=small},
  before skip = 8pt,
  after skip  = 8pt
}

% --- Intuition box (purple) — for plain-language explanations, analogies ---
\newtcolorbox{intuition}[1][]{
  enhanced,
  breakable,
  arc         = 2mm,
  colback     = boxPurpleBg,
  colframe    = boxPurple!50,
  title       = #1,
  fonttitle   = \bfseries\color{white},
  coltitle    = white,
  colbacktitle= boxPurple,
  attach boxed title to top left = {yshift=-2mm, xshift=2mm},
  boxed title style = {arc=2mm, size=small},
  before skip = 8pt,
  after skip  = 8pt
}

% ---------- Beamer theme ----------
\usetheme{Madrid}
\usecolortheme{default}

% --- Custom beamer colors (blue accent matching explanation boxes) ---
\setbeamercolor{structure}{fg=boxBlue}
\setbeamercolor{palette primary}{fg=white, bg=boxBlue}
\setbeamercolor{palette secondary}{fg=white, bg=boxBlue!80}
\setbeamercolor{palette tertiary}{fg=white, bg=boxBlue!60}
\setbeamercolor{block title}{fg=white, bg=boxBlue}
\setbeamercolor{block body}{fg=black, bg=boxBlueBg}

% Remove navigation symbols for clean look
\setbeamertemplate{navigation symbols}{}

% Note: Frame title prefix removed — \addtobeamertemplate conflicts with beamer internals

% --- Outline (Table of Contents) styling ---
\setbeamerfont{section in toc}{family=\headingfont, series=\mdseries}
\setbeamerfont{subsection in toc}{family=\headingfont, series=\mdseries}
\setbeamertemplate{section in toc}{%
  \leavevmode\leftskip=1.5em
  \textcolor{boxBlue}{\textbf{\large\inserttocsectionnumber.}}%
  \quad\inserttocsection\par
  \vspace{0.2em}
}
\setbeamertemplate{subsection in toc}{%
  \leavevmode\leftskip=4em
  \textcolor{boxBlue!60}{\small$\bullet$}%
  \quad\inserttocsubsection\par
  \vspace{0.1em}
}

% Section header slides — Thai in white, English in smaller light blue below
\AtBeginSection{
  \begin{frame}[plain]{}
    \vspace*{2cm}
    \begin{beamercolorbox}[sep=12pt, center, rounded=true, shadow=true]{palette primary}
      \IfSubStr{\secname}{(}{%
        \StrBefore{\secname}{(}[\sectionthai]%
        \StrBetween[1,1]{\secname}{(}{)}[\sectioneng]%
        \begin{tabular}{@{}c@{}}
          {\usebeamerfont{title}\sectionthai} \\[0.2em]
          {\usebeamerfont{subtitle}\color{boxBlue!60}\sectioneng}
        \end{tabular}%
      }{%
        {\usebeamerfont{title}\secname\par}%
      }
    \end{beamercolorbox}
    \vfill
  \end{frame}
}

% Subsection header slides — smaller, centered
\AtBeginSubsection{
  \begin{frame}[plain]{}
    \vfill
    \centering
    {\usebeamerfont{section title}\insertsubsectionhead\par}
    \vfill
  \end{frame}
}

% Minimal footer: page number only, right-aligned
\setbeamertemplate{footline}{
  \hfill\usebeamerfont{page number in head/foot}\insertframenumber{} / \inserttotalframenumber\hspace*{1em}\vspace*{0.5ex}
}

% ---------- Spacing & readability ----------
\linespread{1.2}

% ---------- Useful commands ----------

% Highlight important text in blue
\newcommand{\highlight}[1]{\textcolor{boxBlue}{\textbf{#1}}}

% Math set notation helpers
\newcommand{\R}{\mathbb{R}}
\newcommand{\N}{\mathbb{N}}
\newcommand{\Z}{\mathbb{Z}}
\newcommand{\Q}{\mathbb{Q}}

% ============================================================================
%  End of preamble
% ============================================================================


% --- Pseudo code box (TH Sarabun New — main font, green theme) ---
\newtcolorbox{pseudobox}[1][]{
  enhanced,
  breakable,
  sharp corners,
  colback  = boxGreen!6,
  colframe = boxGreen,
  leftrule = 3mm,
  title    = #1,
  fonttitle     = \bfseries\color{white},
  coltitle      = white,
  colbacktitle  = boxGreen,
  attach boxed title to top left = {yshift=-2mm, xshift=2mm},
  boxed title style = {sharp corners, size=small},
  before skip = 8pt,
  after skip  = 8pt
}

\title[ขั้นตอนวิธี]{ขั้นตอนวิธี}
\subtitle{Algorithms — การติดตาม Pseudo Code และกลยุทธ์การทำโจทย์}
\author[None W. (Yuan)]{None W. (Yuan)}
\date{}

\begin{document}

\begin{frame}[plain]\titlepage\end{frame}
\begin{frame}{สารบัญ (Outline)}\tableofcontents\end{frame}

\begin{frame}{ก่อนเรียน (Prerequisites)}
  \begin{explanation}[สิ่งที่ควรรู้ก่อนเรียน]
    \begin{itemize}
      \item \textbf{การโปรแกรมเบื้องต้น} — if/else, while loop, ตัวแปร, ตัวดำเนินการ
      \item \textbf{คณิตศาสตร์พื้นฐาน} — การบวก ลบ คูณ หาร, การเปรียบเทียบ, มอดุโล
      \item \textbf{ตรรกศาสตร์} — and, or, not, การให้เหตุผล
    \end{itemize}
  \end{explanation}

  \begin{intuition}[ขั้นตอนวิธีคืออะไร?]
    \textbf{ขั้นตอนวิธี (Algorithm)} คือลำดับขั้นตอนการทำงานที่เป็นระบบ
    — ในข้อสอบ สอวน. จะเขียนในรูปแบบ \highlight{Pseudo Code ภาษาไทย}
    โดยใช้เลขกำกับขั้นตอนแบบ 1., 1.1, 2., 3.1.1
  \end{intuition}
\end{frame}

% ===========================================================================
%  SECTION 1: Pseudo Code คืออะไร
% ===========================================================================
\section{Pseudo Code คืออะไร?}

\begin{frame}{Pseudo Code ในข้อสอบ สอวน.}
  \begin{explanation}[รูปแบบ]
    \begin{itemize}
      \item เขียนเป็น \textbf{ภาษาไทย} ใช้เลขกำกับขั้นตอน: 1., 2., 3.
      \item ขั้นตอนย่อยใช้เลขต่อท้าย: 1.1, 1.2, 3.1.1
      \item การเยื้อง (indent) แสดงการอยู่ภายใต้เงื่อนไขหรือการวนซ้ำ
    \end{itemize}
  \end{explanation}

  \begin{columns}[T]
    \column{0.5\textwidth}
      \begin{pseudobox}[ตัวอย่าง: if-else]
        \textbf{1.)} ถ้า $x > 10$ \\
        .\qquad \textbf{1.1)} พิมพ์ $x$ \\
        \textbf{2.)} มิฉะนั้น \\
        .\qquad \textbf{2.1)} พิมพ์ $x+5$
      \end{pseudobox}
    \column{0.5\textwidth}
      \begin{exambox}[เทียบเท่า Python]
        \texttt{if x > 10:}

        \qquad \texttt{print(x)}

        \texttt{else:}

        \qquad \texttt{print(x+5)}
      \end{exambox}
  \end{columns}
\end{frame}

\begin{frame}{การแปลง Python $\leftrightarrow$ Pseudo Code}
  \begin{explanation}
    \begin{tabular}{l|l}
      \textbf{Python} & \textbf{Pseudo Code (ไทย)} \\
      \hline
      \texttt{if} เงื่อนไข: & ถ้า เงื่อนไข แล้ว \\
      \texttt{elif} เงื่อนไข: & มิฉะนั้น ถ้า เงื่อนไข แล้ว \\
      \texttt{else}: & มิฉะนั้น \\
      \texttt{while} เงื่อนไข: & ทำซ้ำเมื่อ เงื่อนไข \\
      \texttt{x = 5} & ให้ $x$ มีค่าเท่ากับ $5$ \\
      \texttt{print(x)} & พิมพ์ $x$ / แสดงผล $x$ \\
      \texttt{break} & ออกจากการวนซ้ำ \\
      \texttt{continue} & ข้ามไปทำรอบถัดไป \\
      รับค่า (\texttt{input}) & รับค่า / รับข้อมูลเข้า \\
      \texttt{n \% m} & เศษที่เหลือจากการหาร $n$ ด้วย $m$ \\
      \texttt{n // m} & ผลหารจำนวนเต็มของ $n$ หารด้วย $m$ \\
    \end{tabular}
  \end{explanation}
\end{frame}

% ===========================================================================
%  SECTION 2: วิธีติดตาม Pseudo Code
% ===========================================================================
\section{วิธีติดตาม Pseudo Code}

\begin{frame}{เทคนิคการติดตาม: ตารางตัวแปร (Variable Table)}
  \begin{explanation}[วิธีพื้นฐานที่สุด — ใช้ได้กับทุกโจทย์]
    \textbf{สร้างตาราง} 1 แถวต่อ 1 การเปลี่ยนแปลงค่า —
    เขียนค่าของทุกตัวแปรหลังจากทำแต่ละขั้นตอน
  \end{explanation}

  \begin{columns}[T]
    \column{0.4\textwidth}
      \begin{pseudobox}[Pseudo Code]
        \textbf{1.)} ให้ $a = 5$ \\
        \textbf{2.)} ให้ $b = 3$ \\
        \textbf{3.)} ให้ $a = a + b$ \\
        \textbf{4.)} ให้ $b = a - b$ \\
        \textbf{5.)} พิมพ์ $a, b$
      \end{pseudobox}
    \column{0.6\textwidth}
      \begin{solbox}
        \begin{tabular}{c|c|c}
          \textbf{ขั้นตอน} & \textbf{a} & \textbf{b} \\
          \hline
          เริ่ม  & -- & -- \\
          1     & 5  & -- \\
          2     & 5  & 3  \\
          3     & 8  & 3  \\
          4     & 8  & 5  \\
          5     & \multicolumn{2}{c}{\highlight{พิมพ์ 8, 5}} \\
        \end{tabular}
      \end{solbox}
  \end{columns}
\end{frame}

\begin{frame}{เทคนิค: นับรอบใน While Loop}
  \begin{columns}[T]
    \column{0.4\textwidth}
      \begin{pseudobox}[Pseudo Code]
        \textbf{1.)} ให้ $i = 1$ \\
        \textbf{2.)} ให้ $s = 0$ \\
        \textbf{3.)} ทำซ้ำเมื่อ $i \leq 4$ \\
        \qquad \textbf{3.1)} $s = s + i$ \\
        \qquad \textbf{3.2)} $i = i + 1$ \\
        \textbf{4.)} พิมพ์ $s$
      \end{pseudobox}
    \column{0.6\textwidth}
      \begin{solbox}
        \begin{tabular}{c|c|c|c}
          \textbf{รอบ} & \textbf{i} & \textbf{i≤4?} & \textbf{s} \\
          \hline
          เริ่ม & 1 & -- & 0 \\
          1 & 1 & True & 1 \\
          & 2 & -- & -- \\
          2 & 2 & True & 3 \\
          & 3 & -- & -- \\
          3 & 3 & True & 6 \\
          & 4 & -- & -- \\
          4 & 4 & True & 10 \\
          & 5 & False & -- \\
        \end{tabular}

        \highlight{พิมพ์ 10}
      \end{solbox}
  \end{columns}
\end{frame}

% ===========================================================================
%  SECTION 3: ประเภทคำถามและกลยุทธ์
% ===========================================================================
\section{ประเภทคำถามและกลยุทธ์}

\begin{frame}{ประเภทที่ 1: ติดตามหาผลลัพธ์ (Trace \& Find Output)}
  \begin{explanation}[คำถาม: ``ข้อใดให้ผลลัพธ์ตรงกับขั้นตอนวิธีต่อไปนี้'']
    \textbf{กลยุทธ์}: สร้างตารางตัวแปร ติดตามทีละขั้นตอน แล้วเทียบกับตัวเลือก
  \end{explanation}

  \begin{columns}[T]
    \column{0.45\textwidth}
      \begin{pseudobox}[Pseudo Code]
        \textbf{1.)} ให้ $x = 12$ \\
        \textbf{2.)} ให้ $y = 7$ \\
        \textbf{3.)} ถ้า $x > y$ แล้ว \\
        .\qquad \textbf{3.1)} $z = x - y$ \\
        \textbf{4.)} มิฉะนั้น \\
        .\qquad \textbf{4.1)} $z = y - x$ \\
        \textbf{5.)} พิมพ์ $z$
      \end{pseudobox}
    \column{0.55\textwidth}
      \begin{solbox}[ติดตาม]
        $x=12, y=7$: $x>y$ เป็นจริง $\rightarrow$ ทำขั้น 3.1

        $z = 12-7 = 5$

        \highlight{พิมพ์ 5}
      \end{solbox}
  \end{columns}
\end{frame}

\begin{frame}{ประเภทที่ 2: นับจำนวนการทำงาน (Count Operations)}
  \begin{explanation}[คำถาม: ``ในขั้นตอนที่ X มีการเปรียบเทียบกี่ครั้ง'']
    \textbf{กลยุทธ์}: นับเฉพาะการเปรียบเทียบที่เงื่อนไขในขั้นตอนนั้น —
    เงื่อนไข while ถูกเช็ค 1 ครั้งต่อรอบ (+1 ครั้งสุดท้ายที่ออกจากวน)
  \end{explanation}

  \begin{columns}[T]
    \column{0.45\textwidth}
      \begin{pseudobox}[Pseudo Code]
        \textbf{1.)} ให้ $n = 10$ \\
        \textbf{2.)} ให้ $c = 0$ \\
        \textbf{3.)} ทำซ้ำเมื่อ $n > 0$ \\
        .\qquad \textbf{3.1)} $n = n - 2$ \\
        .\qquad \textbf{3.2)} $c = c + 1$ \\
        \textbf{4.)} พิมพ์ $c$
      \end{pseudobox}
    \column{0.55\textwidth}
      \begin{solbox}[นับการเปรียบเทียบ \texttt{n>0} ในขั้น 3]
        รอบ 1: $n=10 > 0$ $\checkmark$ $\rightarrow$ $n=8, c=1$

        รอบ 2: $n=8 > 0$ $\checkmark$ $\rightarrow$ $n=6, c=2$

        รอบ 3: $n=6 > 0$ $\checkmark$ $\rightarrow$ $n=4, c=3$

        รอบ 4: $n=4 > 0$ $\checkmark$ $\rightarrow$ $n=2, c=4$

        รอบ 5: $n=2 > 0$ $\checkmark$ $\rightarrow$ $n=0, c=5$

        รอบ 6: $n=0 > 0$ ไม่จริง $\rightarrow$ จบ

        \highlight{เปรียบเทียบ $n>0$ ทั้งหมด 6 ครั้ง}
      \end{solbox}
  \end{columns}
\end{frame}

\begin{frame}{ประเภทที่ 3: กำหนด Input ให้ หา Output}
  \begin{explanation}[คำถาม: ``ถ้า number เป็น X จะแสดงค่าตรงกับข้อใด'']
    \textbf{กลยุทธ์}: แทนค่าที่กำหนดลงในตัวแปรเริ่มต้น แล้วติดตามตามปกติ
  \end{explanation}

  \begin{columns}[T]
    \column{0.567\textwidth}
      \begin{pseudobox}[Pseudo Code]
        \textbf{1.)} รับค่า $n$ \\
        \textbf{2.)} ถ้า $n \bmod 3 = 0$ และ $n \bmod 5 = 0$ แล้ว \\
        .\qquad \textbf{2.1)} พิมพ์ ``A'' \\
        \textbf{3.)} มิฉะนั้น ถ้า $n \bmod 3 = 0$ หรือ $n \bmod 5 = 0$ แล้ว \\
        .\qquad \textbf{3.1)} พิมพ์ ``B'' \\
        \textbf{4.)} มิฉะนั้น \\
        .\qquad \textbf{4.1)} พิมพ์ ``C''
      \end{pseudobox}
    \column{0.49\textwidth}
      \begin{solbox}[เมื่อ n เป็น 50]
        $50 \bmod 3 = 2 \neq 0$ $\rightarrow$ เงื่อนไข 2 เป็นเท็จ

        $50 \bmod 5 = 0$ $\rightarrow$ เงื่อนไข 3 เป็นจริง

        \highlight{พิมพ์ ``B''}
      \end{solbox}

      \begin{exambox}[เมื่อ n เป็น 15]
        $15 \bmod 3 = 0$ และ $15 \bmod 5 = 0$ $\rightarrow$ เงื่อนไข 2 เป็นจริง
        \highlight{พิมพ์ ``A''}
      \end{exambox}
  \end{columns}
\end{frame}

\begin{frame}{ประเภทที่ 4: นับจำนวนครั้งที่ขั้นตอนถูกทำ (Iteration Count)}
  \begin{explanation}[คำถาม: ``ขั้นตอน X ถูกประมวลผลกี่ครั้ง'']
    \textbf{กลยุทธ์}: ดูว่า loop วนกี่รอบ — แต่ละรอบทำขั้นตอนนั้นกี่ครั้ง
    (nested loop: คูณจำนวนรอบ)
  \end{explanation}

  \begin{columns}[T]
    \column{0.45\textwidth}
      \begin{pseudobox}
        \textbf{1.)} ให้ $c = 0$ และ ให้ $i = 1$ \\
        \textbf{2.)} ทำซ้ำเมื่อ $i \leq 3$ \\
        .\qquad \textbf{2.1)} ให้ $j = 1$ \\
        .\qquad \textbf{2.2)} ทำซ้ำเมื่อ $j \leq 2$ \\
        .\qquad\quad \textbf{2.2.1)} $c = c + 1$ \\
        .\qquad\quad \textbf{2.2.2)} $j = j + 1$ \\
        .\qquad \textbf{2.3)} $i = i + 1$ \\
        \textbf{3.)} พิมพ์ $c$
      \end{pseudobox}
    \column{0.55\textwidth}
      \begin{solbox}[ขั้น 2.2.1 ถูกทำกี่ครั้ง?]
        Loop นอก ($i$) วน 3 รอบ

        Loop ใน ($j$) วน 2 รอบต่อ 1 รอบของ loop นอก

        ขั้น 3.2.1 ($c = c+1$) อยู่ใน loop ใน

        จำนวนครั้ง $= 3 \times 2 =$ \highlight{6 ครั้ง}

        \medskip
        (แต่ละครั้ง $c$ เพิ่มทีละ 1 $\rightarrow$ พิมพ์ $c = 6$)
      \end{solbox}
  \end{columns}
\end{frame}

\begin{frame}{ประเภทที่ 5: หาจุดผิดพลาด (Find Error/Bug)}
  \begin{explanation}[คำถาม: ``ปรากฏว่าอัลกอริทึมไม่ได้ให้ผลที่ถูกต้อง บรรทัดใดมีส่วนที่ไม่ถูกต้อง'']
    \textbf{กลยุทธ์}: ตรวจสอบทีละจุด — เงื่อนไข, การกำหนดค่าเริ่มต้น, ขอบเขต, การออกจากวน
  \end{explanation}

  \begin{columns}[T]
    \column{0.45\textwidth}
      \begin{pseudobox}[Pseudo Code ในการหาผลรวม 1-4]
        \textbf{1.)} ให้ $i = 1$ \\
        \textbf{2.)} ให้ $s = 0$ \\
        \textbf{3.)} ทำซ้ำเมื่อ $i < 5$ \\
        .\qquad \textbf{3.1)} $s = s + i$ \\
        \textbf{4.)} พิมพ์ $s$
      \end{pseudobox}
    \column{0.55\textwidth}
      \begin{solbox}[อะไรผิด?]
        วนรอบ: $i=1,2,3,4$ $\rightarrow$ $s=1+2+3+4=10$

        แต่ $i$ \textbf{ไม่เคยเปลี่ยนค่า} ใน loop!

        ขั้น 3.1 ไม่มีการเพิ่ม $i$ \quad $\rightarrow$ \highlight{loop อนันต์!}

        \textbf{ควรเพิ่ม:} 3.2) $i = i + 1$
      \end{solbox}
  \end{columns}
\end{frame}

\begin{frame}{ประเภทที่ 6: การจัดการ Array/String}
  \begin{explanation}[คำถาม: ``ถ้า data1 เป็น ... data2 เป็น ... ผลลัพธ์ในขั้นตอนที่ X คืออะไร'']
    \textbf{กลยุทธ์}: ติดตามค่าแต่ละ array แยกกัน — เขียน array ทั้งสองแถว อัปเดตเมื่อมีการเปลี่ยนแปลง
  \end{explanation}

  \begin{columns}[T]
    \column{0.5\textwidth}
      \begin{pseudobox}
        \textbf{1.)} ให้ $A = [2, 5, 8]$ \\
        \textbf{2.)} ให้ $B = [1, 3, 6]$ \\
        \textbf{3.)} ให้ $i = 1$ \\
        \textbf{4.)} ทำซ้ำเมื่อ $i \leq 3$ \\
        .\qquad \textbf{4.1)} ถ้า $i$ เป็นเลขคี่ แล้ว \\
        .\qquad\quad \textbf{4.1.1)} สลับ $A[i]$ กับ $B[i]$ \\
        .\qquad \textbf{4.2)} $i = i + 1$ \\
        \textbf{5.)} แสดงผล $A$
      \end{pseudobox}
    \column{0.5\textwidth}
      \begin{solbox}[ติดตาม]
        \textbf{เริ่ม:} $A=[2,5,8]$, $B=[1,3,6]$

        $i=1$ (คี่): สลับ $A[1],B[1]$

        $\rightarrow$ $A=[1,5,8]$, $B=[2,3,6]$

        $i=2$ (คู่): ไม่สลับ

        $i=3$ (คี่): สลับ $A[3],B[3]$

        $\rightarrow$ $A=[1,5,6]$, $B=[2,3,8]$

        \highlight{$A = [1, 5, 6]$}
      \end{solbox}
  \end{columns}
\end{frame}

\begin{frame}{ประเภทที่ 7: ความเหมือนเชิงตรรกะ (Logical Equivalence)}
  \begin{explanation}[คำถาม: ``ข้อใดให้ผลลัพธ์เหมือนขั้นตอนวิธีต่อไปนี้'']
    \textbf{กลยุทธ์}: แปลง pseudo code เป็นนิพจน์ตรรกะ หรือทดสอบด้วย input ง่าย ๆ
  \end{explanation}

  \begin{columns}[T]
    \column{0.4\textwidth}
      \begin{pseudobox}[Pseudo Code]
        \textbf{1.)} รับค่า $x$ \\
        \textbf{2.)} ถ้า $x > 10$ แล้ว \\
        .\qquad \textbf{2.1)} ถ้า $x < 20$ แล้ว \\
        .\qquad\quad \textbf{2.1.1)} พิมพ์ ``Yes'' \\
        .\qquad \textbf{2.2)} มิฉะนั้น พิมพ์ ``No'' \\
        \textbf{3.)} มิฉะนั้น พิมพ์ ``No''
      \end{pseudobox}
    \column{0.6\textwidth}
      \begin{solbox}[สมมูลกับ]
        พิมพ์ ``Yes'' เมื่อ \highlight{$x > 10$ และ $x < 20$}
        (คือ $10 < x < 20$)

        นอกนั้นพิมพ์ ``No''

        \medskip
        \textbf{เทียบเท่า:}

        \texttt{if x > 10 and x < 20:}

        \qquad \texttt{print("Yes")}

        \texttt{else: print("No")}
      \end{solbox}
  \end{columns}
\end{frame}

% ===========================================================================
%  SECTION 4: Mod ใน Pseudo Code
% ===========================================================================
\section{โจทย์ Mod ใน Pseudo Code}

\begin{frame}{โจทย์ Mod: การวนรอบและวันในสัปดาห์}
  \begin{explanation}[คำถามแนว Mod]
    โจทย์ mod มักถามเรื่อง \textbf{การวนรอบ} — วันในสัปดาห์, วงล้อหมุน, การเลื่อนตัวอักษร\\
    \textbf{กลยุทธ์}: หาความยาวรอบ $\rightarrow$ หาเศษจากการหาร $\rightarrow$ ตอบ
  \end{explanation}

  \begin{columns}[T]
    \column{0.6\textwidth}
      \begin{pseudobox}[ตัวอย่าง: วันในสัปดาห์]
        \textbf{1.)} รับค่า $d$ (จำนวนวัน) \\
        \textbf{2.)} ให้ $r =$ เศษที่เหลือจากการหาร $d$ ด้วย $7$ \\
        \textbf{3.)} ถ้า $r = 0$ แล้ว พิมพ์ ``อาทิตย์'' \\
        \textbf{4.)} มิฉะนั้น ถ้า $r = 1$ แล้ว พิมพ์ ``จันทร์'' \\
        \quad\vdots \\
        \textbf{9.)} มิฉะนั้น พิมพ์ ``เสาร์''
      \end{pseudobox}

  \column{0.4\textwidth}
    \begin{intuition}[หลักการ]
      mod 7 = ดูว่าเหลือเศษกี่วันหลังจากครบสัปดาห์ —
      \textbf{รอบละ 7 วัน} วนกลับมาที่เดิมเสมอ
    \end{intuition}
  \end{columns}
\end{frame}

% ===========================================================================
%  SECTION 5: ตัวอย่างขั้นสูง
% ===========================================================================
\section{ตัวอย่างขั้นสูง (Advanced Tracing)}

\begin{frame}{ตัวอย่าง: การวนซ้ำซับซ้อน}
  \begin{exambox}
    จากขั้นตอนวิธีต่อไปนี้ จะพิมพ์ค่าใดออกมา?
  \end{exambox}

  \begin{columns}[T]
    \column{0.45\textwidth}
      \begin{pseudobox}
        \textbf{1.)} ให้ $n = 25$ \\
        \textbf{2.)} ให้ $c = 0$ \\
        \textbf{3.)} ทำซ้ำเมื่อ $n > 1$ \\
        .\qquad \textbf{3.1)} ถ้า $n \bmod 2 = 0$ แล้ว \\
        .\qquad\quad \textbf{3.1.1)} $n = n // 2$ \\
        .\qquad \textbf{3.2)} มิฉะนั้น \\
        .\qquad\quad \textbf{3.2.1)} $n = n - 1$ \\
        .\qquad \textbf{3.3)} $c = c + 1$ \\
        \textbf{4.)} พิมพ์ $c$
      \end{pseudobox}
    \column{0.55\textwidth}
      \begin{solbox}
        $n=25$ (คี่) $\rightarrow$ $n=24$, $c=1$

        $n=24$ (คู่) $\rightarrow$ $n=12$, $c=2$

        $n=12$ (คู่) $\rightarrow$ $n=6$, $c=3$

        $n=6$ (คู่) $\rightarrow$ $n=3$, $c=4$

        $n=3$ (คี่) $\rightarrow$ $n=2$, $c=5$

        $n=2$ (คู่) $\rightarrow$ $n=1$, $c=6$

        $n=1$: $n > 1$ เป็นเท็จ $\rightarrow$ จบ

        \highlight{พิมพ์ 6}
      \end{solbox}
  \end{columns}
\end{frame}

\begin{frame}{ตัวอย่าง: โจทย์ Fill-in (4 หลัก)}
  \begin{explanation}[กลยุทธ์: อย่าตกใจ — ติดตามด้วยมือได้เสมอ]
    โจทย์ fill-in มักให้กรณีทดสอบขนาดเล็ก
    แม้จะเป็น Dijkstra, DP, หรือ DFS ก็ตาม
  \end{explanation}

  \begin{exambox}[ตัวอย่าง: Dijkstra 4 โหนด 5 เส้น]
    A-B:3, A-C:5, B-C:1, B-D:7, C-D:2 เริ่มที่ A หาระยะสั้นสุดไป D

    \begin{tabular}{c|c|c|c}
      \textbf{ขั้น} & A & B & C \\
      \hline
      เริ่ม & 0 & $\infty$ & $\infty$ \\
      จาก A & 0 & 3 & 5 \\
      จาก B & 0 & 3 & 4 (ผ่าน B) \\
      จาก C & 0 & 3 & 4 \\
    \end{tabular}

    \highlight{A$\to$D = 4+2 = 6 \quad (ตอบ 0006)}
  \end{exambox}
\end{frame}

% ===========================================================================
%  SUMMARY
% ===========================================================================
\section{สรุป}

\begin{frame}{สรุปกลยุทธ์การทำโจทย์ Pseudo Code}
  \begin{explanation}[7 ประเภทคำถาม + กลยุทธ์]
    \begin{columns}[T]
      \column{0.5\textwidth}
        \begin{enumerate}
          \item \textbf{ติดตามหาผลลัพธ์} — ตารางตัวแปร
          \item \textbf{นับจำนวนการทำงาน} — นับเฉพาะในเงื่อนไข
          \item \textbf{input $\to$ output} — แทนค่าแล้วติดตาม
          \item \textbf{นับ iteration} — loop นอก $\times$ loop ใน
        \end{enumerate}
      \column{0.5\textwidth}
        \begin{enumerate}
          \setcounter{enumi}{4}
          \item \textbf{หาจุดผิดพลาด} — infinite loop, เงื่อนไขผิด
          \item \textbf{array/string} — ติดตามแยกทีละ array
          \item \textbf{สมมูลเชิงตรรกะ} — แปลงเป็น and/or
        \end{enumerate}
    \end{columns}

    \highlight{ทุกโจทย์แก้ได้ด้วย: ติดตามทีละขั้น + ตารางตัวแปร!}
  \end{explanation}
\end{frame}

\end{document}
